\documentclass[11pt]{article}
\usepackage[T1]{fontenc}
\usepackage{lmodern}
\usepackage[margin=1in]{geometry}
\usepackage{amsmath,amssymb,microtype,hyperref}
\title{Gaussian and Eisenstein Parity Tables\\\large Fifty-six proved specializations at weights two through eight}
\author{Companion to \emph{Mixed-Point Elimination and Certified Parity Reductions}}
\date{October 9, 2026}
\begin{document}
\maketitle
The convention is $F_{a,b}(z)=\operatorname{Li}_{a,b}(z,1)$ with decreasing nested indices.
Write $\rho=e^{2\pi i/3}$,
$B_{2r}=\beta(2r)$, $C_{2r}=\Im\operatorname{Li}_{2r}(\rho)$,
$\ell_2=\log2$ and $\ell_3=\log3$.
The projection is imaginary at even weight and real at odd weight.
These expressions are consequences of the explicit parity theorem in the accompanying article,
not independent PSLQ conjectures. They are generated by exact symbolic substitution.
The Gaussian block comes first, followed by the Eisenstein block; each block is ordered by weight and first index.
\input{parity_tables.tex}
\end{document}
